\section{Conclusion}

This paper introduced Domphy and the Compositional Atomic Reactive
Mapping (CARM) model, a flat compositional architecture for declarative
user interfaces. By replacing hierarchical component recursion with a
single-pass structural reification process, the model establishes a
persistent shell in which reactive bindings are constructed once and
maintained throughout the application lifecycle.

The proposed atomic mapping mechanism enables bounded, slot-level
updates that are independent of global tree size. Under bounded
dependency assumptions, the update path achieves constant-time
complexity per affected attribute. Empirical observations indicate
reduced structural depth and simplified source definitions compared to
conventional component-based frameworks.

While the model introduces additional memory overhead due to persistent
reactive bindings, it offers deterministic runtime behavior and predictable
latency characteristics. These properties make the architecture suitable
for interaction-intensive environments requiring stable frame timing and
fine-grained updates.

Future work includes formal verification of dependency constraints,
large-scale benchmarking under heterogeneous workloads, and exploration
of static analysis techniques to optimize memory residency.
